%*******************************************************
%  classicthesis-config-pt.tex
%  Configuração ClassicThesis adaptada para português (PT-BR)
%  Baseada no classicthesis-config.tex de André Miede
%*******************************************************

% ── Cores (deve vir antes de tikz ou qualquer pacote que carregue xcolor) ───
\PassOptionsToPackage{dvipsnames,svgnames,x11names}{xcolor}
\usepackage{xcolor}

% ── Codificação ──────────────────────────────────────────────
\PassOptionsToPackage{utf8}{inputenc}
  \usepackage{inputenc}
\PassOptionsToPackage{T1}{fontenc}
  \usepackage{fontenc}

% ── ClassicThesis (Palatino) ─────────────────────────────────
\PassOptionsToPackage{
  drafting=false,
  tocaligned=false,
  dottedtoc=false,
  eulerchapternumbers=false,
  floatperchapter=true,
  eulermath=false,
  beramono=false,
  palatino=true,
  style=classicthesis
}{classicthesis}

% ── Identificação do documento ───────────────────────────────
\newcommand{\myTitle}{Teoria de Redes de Petri Hierárquica de Sinais\xspace}
\newcommand{\mySubtitle}{Um Arcabouço Formal para Controle Regulatório
  Multi-Escala em Sistemas Celulares\xspace}
\newcommand{\myDegree}{Professor Titular\xspace}
\newcommand{\myName}{Eugênio Simão\xspace}
\newcommand{\myUni}{Universidade Federal de Santa Catarina\xspace}
\newcommand{\myFaculty}{Departamento de Engenharia de Computação\xspace}
\newcommand{\myDepartment}{Campus Araranguá\xspace}
\newcommand{\myLocation}{Araranguá, Santa Catarina\xspace}
\newcommand{\myTime}{Novembro de 2025\xspace}

% ── Abreviaturas ─────────────────────────────────────────────
\providecommand{\ie}{i.\,e.}
\providecommand{\Ie}{I.\,e.}
\providecommand{\eg}{e.\,g.}
\providecommand{\Eg}{E.\,g.}
\newcommand{\cf}{cf.\xspace}

% ── Idioma ───────────────────────────────────────────────────
\PassOptionsToPackage{brazil,english}{babel}
  \usepackage{babel}
\usepackage{csquotes}

% ── Matematica ───────────────────────────────────────────────
\PassOptionsToPackage{fleqn}{amsmath}
  \usepackage{amsmath}
\usepackage{mathtools}   % estende amsmath: dcases, \mathclap, \shortintertext
\usepackage{amssymb}
\usepackage{amsfonts}
\usepackage{amsthm}

% ── Gráficos ─────────────────────────────────────────────────
\usepackage{graphicx}
\usepackage{float}          % [H] placement specifier
\usepackage{scrhack}
\usepackage{xspace}

% ── Tabelas ──────────────────────────────────────────────────
\usepackage{tabularx}
  \setlength{\extrarowheight}{3pt}
\usepackage{booktabs}
\usepackage{longtable}
\usepackage{array}
\usepackage{multirow}
\newcommand{\tableheadline}[1]{\multicolumn{1}{l}{\spacedlowsmallcaps{#1}}}
\newcommand{\myfloatalign}{\centering}
\usepackage{subfig}
\usepackage{wrapfig}        % text-wrapped figures

% ── Algoritmos ───────────────────────────────────────────────
\usepackage{algorithm}
\usepackage{algpseudocode}

% ── Química ──────────────────────────────────────────────────
\usepackage[version=4]{mhchem}

% ── TikZ (Redes de Petri) ────────────────────────────────────
\usepackage{tikz}
\usetikzlibrary{arrows.meta, positioning, petri, shapes, calc}

% ── Símbolos (pifont para \ding, etc.) ───────────────────────
\usepackage{pifont}

% ── Bibliografia (biblatex/biber) ─────────────────────────────
\PassOptionsToPackage{%
  backend=biber,
  citestyle=abnt,
  bibstyle=numeric,
  sorting=nyt,
  natbib=true,
  hyperref=true
}{biblatex}
\usepackage{biblatex}

% ── Referência à última página (ficha catalográfica) ────────
\usepackage{lastpage}

% ── Microtype ────────────────────────────────────────────────
\usepackage{microtype}

% ── Ambientes matemáticos ────────────────────────────────────
\theoremstyle{definition}
\newtheorem{definition}{Definição}[chapter]
\newtheorem{principle}[definition]{Princípio}
\newtheorem{example}[definition]{Exemplo}
\theoremstyle{plain}
\newtheorem{theorem}[definition]{Teorema}
\newtheorem{lemma}[definition]{Lema}
\newtheorem{corollary}[definition]{Corolário}
\newtheorem{proposition}[definition]{Proposição}
\theoremstyle{remark}
\newtheorem{remark}[definition]{Observação}

% ── Símbolo QED ──────────────────────────────────────────────
\renewcommand{\qedsymbol}{$\blacksquare$}

% ── Operadores matemáticos (predicados do formalismo SHPN) ───
% Usados nas provas formais; \DeclareMathOperator produz espaçamento correto
\DeclareMathOperator{\Enabled}{Enabled}
\DeclareMathOperator{\Fire}{Fire}
\DeclareMathOperator{\NEnabled}{NormalEnabled}
\DeclareMathOperator{\TEnabled}{TestEnabled}
\DeclareMathOperator{\SEnabled}{SignalEnabled}
\DeclareMathOperator{\Preemption}{PreemptionCheck}

% ── Notação de pré/pós-conjunto (Redes de Petri) ─────────────
% Pré-conjunto normal:        \npre{t}  →  {}^\bullet t
\newcommand{\npre}[1]{{}^\bullet{#1}}
% Pós-conjunto normal:        \npost{t} →  t^\bullet
\newcommand{\npost}[1]{{#1}^\bullet}
% Pré-conjunto de sinal:      \sigpre{t} → {}^\bullet_s t
\newcommand{\sigpre}[1]{{}^\bullet_s{#1}}
% Pós-conjunto de sinal:      \sigpost{t} → t^\bullet_s
\newcommand{\sigpost}[1]{{#1}^\bullet_s}

% ── Notação de alcançabilidade e semântica de compromisso ─────
% Conjunto de alcançabilidade:   \Reach{M_0}  →  \mathcal{R}(M_0)
\newcommand{\Reach}[1]{\mathcal{R}(#1)}
% Espaço de estados:             \StateSpace  →  \mathcal{M}
\newcommand{\StateSpace}{\mathcal{M}}
% Marcação de compromisso:       \Mcommit{p}  →  M_{\mathrm{commit}}(p)
\newcommand{\Mcommit}[1]{M_{\mathrm{commit}}(#1)}
% Passo de disparo:              \firestep{M}{t}{M'}  →  M \xrightarrow{t} M'
\newcommand{\firestep}[3]{#1 \xrightarrow{#2} #3}
% Alcançabilidade multi-passo:   \reaches{M}{M'}     →  M \xrightarrow{*} M'
\newcommand{\reaches}[2]{#1 \xrightarrow{*} #2}

% ── Espaçamento entre parágrafos (≈ novos blocos de ideia) ───
\setlength{\parskip}{3pt plus 1pt minus 0pt}

% ── Pacote classicthesis ─────────────────────────────────────
\PassOptionsToPackage{hyperfootnotes=false}{hyperref}
\usepackage{classicthesis}

% ── Legenda de figuras/tabelas: alinhamento à extensão (sem recuo) ───────────
\captionsetup{format=plain, justification=justified}

% ── Hyperref ─────────────────────────────────────────────────
\hypersetup{%
  colorlinks=true, linktocpage=true,
  pdfstartpage=3, pdfstartview=FitV,
  breaklinks=true, pageanchor=true,
  pdfpagemode=UseNone,
  plainpages=false, bookmarksnumbered,
  bookmarksopen=true, bookmarksopenlevel=1,
  hypertexnames=true, pdfhighlight=/O,
  urlcolor=CTurl, linkcolor=CTlink, citecolor=CTcitation,
  pdftitle={\myTitle},
  pdfauthor={\textcopyright\ \myName, \myUni},
  pdfsubject={Tese de Concurso para Professor Titular – UFSC},
  pdfkeywords={Redes de Petri, Biologia de Sistemas, Hierarquia de Sinais,
               Bacillus subtilis, UFSC},
  pdfcreator={pdfLaTeX},
  pdfproducer={LaTeX com hyperref e classicthesis}
}

% ── Autoref em português ──────────────────────────────────────
\makeatletter
\@ifpackageloaded{babel}{%
  \addto\extrasbrazil{%
    \renewcommand*{\figureautorefname}{Figura}%
    \renewcommand*{\tableautorefname}{Tabela}%
    \renewcommand*{\partautorefname}{Parte}%
    \renewcommand*{\chapterautorefname}{Capítulo}%
    \renewcommand*{\sectionautorefname}{Seção}%
    \renewcommand*{\subsectionautorefname}{Seção}%
    \renewcommand*{\subsubsectionautorefname}{Seção}%
    \renewcommand*{\equationautorefname}{Eq.}%
  }%
}{\relax}
\makeatother
